\section{讲座定位：为什么 ``训练配方'' 比单点技巧更重要}

\subsection{问题定义：Reasoning 不是单一能力，而是训练链路的系统产物}
本讲最核心的价值不是再讲一个新算法，而是把 ``reasoning 能力'' 拆回工程现实：它由预训练分布、后训练数据、优化目标、验证机制、推理时预算共同决定。也就是说，推理能力不是某个 magical trick 的直接结果，而是一个端到端系统设计问题。

\begin{importantbox}{讲者的核心立场}
\textit{``Open scientific research is why we are here.''} 讲者把过去几年 LLM 进展归因于开放科研生态，并强调当生态变得封闭，研究可复现性、可比较性和可纠错性都会下降。
\end{importantbox}

\begin{knowledgebox}{为什么这一讲要强调 ``recipes''}
模型能力讨论常被简化为参数规模与算力，但工业可复现的 ``recipe'' 其实包含更多隐含变量：
\begin{itemize}
    \item 数据采样与混合策略；
    \item 训练阶段切分（pre-train / post-train / inference-time）；
    \item 评估方式与去污染策略；
    \item 预算约束下的最优折中。
\end{itemize}
\end{knowledgebox}

\begin{figure}[H]
\centering
\includegraphics[width=0.9\textwidth]{frame-01-open-science.jpg}
\caption{视频约 00:03:10：讲者明确把 ``开放研究'' 作为训练配方方法论的前提}
\end{figure}

\subsection{从 ``模型发布'' 到 ``科学发布''}
讲者把开放模型分成三个层级：仅开放权重、开放训练代码、开放数据与评估流程。真正支持科学研究的是第三层，因为只有同时开放数据处理、训练脚本和评估基准，研究社区才能判断改进来自哪里。

\begin{warningbox}{只开权重不等于可复现}
很多工作只发布 checkpoint，导致研究者无法验证提升是来自：
\begin{itemize}
    \item 数据质量改善；
    \item 训练超参数微调；
    \item 优化器与调度修改；
    \item 评估污染或基准差异。
\end{itemize}
当这些变量不可见时，所谓 ``SOTA'' 对科学社区的价值会明显下降。
\end{warningbox}

\subsection{本章小结}
本章给出整场讲座的方法论基线：推理能力要从 ``完整训练配方'' 理解，开放与可复现是该配方可累积迭代的必要条件。


